\chapter{Conclusion and Outlook}\label{ch:conclusion}

\section{Summary of Contributions}
\label{sec:concl:summary}
% State the most important outcomes:
% 1. Introduced a 4-stage pipeline for architecture-aware domain model configuration
% 2. Proved that architecture-aware prompting mathematically outperforms generic prompting
% 3. Demonstrated that a closed feedback loop increases validation/improvement rates significantly
% 4. Showed that the engineering methodology matters more than the choice of LLM

\section{Answers to Research Questions}
\label{sec:concl:rq_answers}
% Copy the research questions from Chapter 1 and provide definitive, data-backed answers
% Each question gets a clear response grounded in the findings from Chapter 7

\section{Limitations}
\label{sec:concl:limitations}
% Software and practical limitations ONLY (experimental validity is in Chapter 8)
% API latency: feedback loop takes several minutes per domain
% Cost: token consumption across multiple LLM calls
% Manual rule engineering: architecture-aware rules were hand-crafted from planner source code
% Single-machine execution: limits the scale of experiments that can be run

\section{Future Work}
\label{sec:concl:future_work}
% Each limitation should spawn an idea for future work
% Automated rule extraction: use a separate LLM agent to read planner C++ source code and extract rules
% Scaling to more planners/domains: temporal, HTN, numeric PDDL
% Online adaptation: real-time feedback during planning, not just offline batch processing
% Fine-tuning LLMs specifically for PDDL configuration tasks
% Integration into planning competition infrastructure
